\documentclass[10pt]{article}
\usepackage[margin=0.75in]{geometry}
\usepackage{amsmath,microtype}
\setlength{\emergencystretch}{3em}
\title{Guide: Where to Look for an Unbounded Orbit}
\author{Collatz proof project}
\date{September 5, 2026}
\begin{document}
\maketitle
\section*{What Lean now proves}
Every hypothetical unbounded positive Collatz orbit has a point from which
its multiplicative coefficient never drops below one. Such starting points
occur arbitrarily late. Conversely, a positive seed whose coefficient
never drops below one must have an unbounded orbit.

This turns the escape problem into one exact task: prove that every positive
seed eventually has a coefficient below one. That task is still open.
Even after solving it, nontrivial cycles would still need to be excluded.

\section*{Why a suitable starting point exists}
The earlier summability result proves that the inverse coefficient
$2^t/3^{j_t}$ tends to zero along any unbounded orbit. Its value at time zero
is one, so its largest value is attained at some finite time.

Start at that maximum. Every later inverse coefficient, measured relative
to the new starting point, is at most one. This is precisely the required
noncontraction condition. The same argument can start after any prescribed
time; it does not say that every late point has the property.

For the converse, suppose such a seed had a bounded orbit. The exact orbit
identity would then bound its ideal correction too. But bounded ideal
correction already forces escape, giving a contradiction.

\section*{A first arithmetic restriction}
A seed with no coefficient contraction must be odd. If it is also two
modulo three, then $(2N-1)/3$ is a smaller positive odd predecessor. Its
first coefficient is $3/2$, so it also has no future contraction.

Therefore a smallest positive seed with this property cannot be two modulo
three. This does not eliminate the remaining residues. Lean also proves
a general rule for transporting noncontraction backward through a finite
prefix whose coefficients all stay at least one.

\section*{What this changes, and what it does not}
The reduction gives complete coverage of hypothetical unbounded positive
orbits by a specific arithmetic class. We still have to prove that this
class has no positive natural seed. Passing a finite prefix test cannot
establish the infinite condition.

The earlier 65-step descent certificate helps when a coefficient
contraction happens within that horizon. It does not prove that every
seed contracts within 65 steps, or at any time. The Collatz goal remains
open; no mathematical priority claim is made here.

\section*{Verification}
The build and 139-declaration standard-foundations audit pass.
{\raggedright Run \texttt{lake build Collatz.NoncontractingTail} from
\texttt{lean/}.\par}
The technical paper gives the exact equivalences and hypotheses.
\end{document}
